\chapter{Sustainability Plan \& Ethical Considerations}\label{ch:sustain}

\section{Scalability}

RecaLL scales differently from a hosted web application. Since each student runs a private copy,
adding users adds no load to any shared server: the work is spread across the users' own machines,
and the only shared resources are the AI providers, which are billed to each student's own key.
Scalability is therefore mainly a question of how much material one student can hold and how fast
one computer can process it.

\begin{itemize}
  \item \textbf{Data volume.} SQLite in Write-Ahead Logging mode allows reads while a write is in
  progress, and the tables used by the review queue and analytics are indexed by workspace, document
  and due date. Large books are processed in stages with checkpoints, so a long PDF can resume after
  an interruption instead of restarting. In the volume test (\Cref{ch:test}) the text of a 930-page textbook
  was extracted in 3.1~s and searched in 0.72~s, and the library of 11 documents and 1{,}070 chunks needs a
  20.8~MB database.
  \item \textbf{AI throughput.} In the study features, provider calls are retried with exponential
  back-off (\texttt{app/llm\_providers/retry.py}) and Gemini calls also pass a priority token-bucket rate
  limiter (\texttt{rate\_limiter.py}), so bursts of note generation do not exceed provider quotas. The
  same policy will also be applied to the \texttt{/api/v1} services as part of the maintenance plan. A student who needs more throughput can move to a faster provider or a local model without
  changing any study data.
  \item \textbf{Towards a hosted version.} The repository already contains a container setup with
  PostgreSQL~16 and pgvector, an Nginx gateway and health checks. Moving the data layer from SQLite to
  PostgreSQL and using semantic embeddings by default would allow a
  multi-user server with semantic search. Workspaces and per-user isolation are already part of the
  \texttt{/api/v1} design, which eases this move.
\end{itemize}

\section{Flexibility / Customization}

\begin{itemize}
  \item \textbf{Choice of AI provider.} Providers sit behind interfaces selected at run time: the study
  features use \texttt{BaseLLMProvider} with a factory, and the \texttt{/api/v1} services use
  \texttt{AIService} with provider adapters. Adding a provider means adding one class per stack;
  students can switch between Gemini, OpenAI, Groq and a local Ollama model at any time.
  \item \textbf{Editable prompts.} The prompts of the study features live in a central registry
  (\texttt{app/prompt\_manager.py}). They can be overridden from the LLM Inspection page; placeholders
  are validated before saving, and any prompt can be reset to its default.
  \item \textbf{Pluggable extractors.} The study path extracts PDFs through a registry of extractors:
  PyMuPDF4LLM is the default and Marker can be added for scanned or handwritten documents.
  \item \textbf{Replaceable scheduler.} Review scheduling sits behind the \texttt{ReviewScheduler}
  interface; the FSRS parameters (desired retention 0.9, maximum interval 36{,}500 days) are
  configuration, and a different algorithm could be plugged in.
  \item \textbf{Personal presentation.} Notes can use ruled, grid or plain paper and a light or dark
  theme, and the interface remembers the student's preferences.
\end{itemize}

\section{Maintainability and Future Enhancements}

\paragraph{Maintenance plan.}
\begin{enumerate}
  \item \textbf{Keep dependencies current.} Pin the Python requirements and review updates monthly, so
  that builds stay reproducible.
  \item \textbf{Unify the two databases.} RecaLL uses \texttt{recall\_saas.db} and a per-user study
  database, and books are mirrored between them. Moving the study routes (books, topics, notes,
  annotations, drills) under \texttt{/api/v1} with the same authentication will let one database serve
  every feature.
  \item \textbf{Documentation.} Keep the README, the SDLC diagrams and this report in step with the code.
\end{enumerate}

\paragraph{Future enhancements.}
\begin{itemize}
  \item Semantic embeddings by default, including from the student's own key, for Knowledge Hub search.
  \item Built-in OCR for scanned PDFs without the optional Marker install.
  \item Smarter scheduling across subjects, using each student's own review history to tune FSRS.
  \item Cloud sync between a student's devices.
  \item macOS and Linux installers, and later a web or mobile companion.
  \item Payments, if a hosted tier is offered.
\end{itemize}

\section{Ethical Considerations}

\begin{itemize}
  \item \textbf{Privacy by design.} Documents, notes, flashcards and study history stay on the student's
  computer. Text is sent to an AI provider only when the student uses an AI feature, and only to the
  provider the student selected; choosing Ollama keeps everything offline.
  \item \textbf{Security of personal keys.} Provider keys are encrypted with AES-256-GCM in the BYOK
  vault and only a short hint is displayed; passwords are stored as PBKDF2-HMAC-SHA256 hashes with
  600{,}000 iterations (\Cref{sec:security}).
  \item \textbf{Honesty about AI output.} AI-generated notes, questions and answers can be wrong.
  Knowledge Hub answers cite the retrieved source passages, and a citation validator removes citations
  that do not match retrieved chunks, so students can check each claim against the original document.
  \item \textbf{Prompt-injection safety.} Retrieved document text is passed to the model as untrusted
  reference data, and the system prompt tells the model not to follow instructions inside it.
  \item \textbf{Learning, not shortcuts.} The design pushes students to explain ideas themselves
  (Socratic drills) and to recall answers before seeing them (flashcards), rather than letting the AI do
  the work for them.
  \item \textbf{Copyright.} Students upload their own course material for personal study; the application
  does not publish or share it.
  \item \textbf{Voluntary participation.} The six students in the user tests (\Cref{sec:uat}) took part
  voluntarily and used the team's own copy of the app; no synthetic or third-party personal data was
  used.
  \item \textbf{Responsible use of AI in this project.} The Antigravity agent generated most of the code
  from requirements and designs the team had prepared beforehand, and was used to review the application
  and find bugs; all output was reviewed and tested by the team. The automated tests were likewise largely
  written with the agent and reviewed by the team (see Appendix~F).
\end{itemize}
